%includo il file che contiene la versione dei documenti
\newcommand{\versioneAnalisiDeiRequisiti}{3.2.0}			
\newcommand{\versioneNormeDiProgetto}{3.2.0}			
\newcommand{\versioneGlossario}{3.2.0}			
\newcommand{\versionePianoDiQualifica}{3.2.0}			
\newcommand{\versionePianoDiProgetto}{3.2.0}	
\newcommand{\versioneStudioDiFattibilita}{1.2.0}
\newcommand{\versioneSpecificaTecnica}{3.2.0}


\newcommand{\Versione}{\versionePianoDiProgetto{}}	%Versione Finale
\newcommand{\Data}{2013-11-28}						%Data di creazione
\newcommand{\DataUltimaModifica}{2013-12-18}
\newcommand{\TipoDocumento}{Piano di Progetto}		%tipo documento

%includo il file header.tex (logo grande in prima pagina piu qualche altra regola)
%questo file contiene impostazioni comuni per tutte i documenti

%definizione packages utilizzati
\documentclass[a4paper]{article}
\usepackage[utf8x]{inputenc}
\usepackage{enumitem}
\usepackage[italian]{babel}
\usepackage{latexsym}
\usepackage{xparse}
\usepackage{float}
\usepackage{subfloat}
\usepackage{subfig}
\usepackage{fancyhdr}
\usepackage{eurofont}
\usepackage{lastpage}
\usepackage{graphicx}
\usepackage{textcomp}
\usepackage{booktabs}
\usepackage{color}
\usepackage{lscape}
\usepackage{hyperref}
%\usepackage[bookmarks=false]{hyperref}
\hypersetup{colorlinks=true, linkcolor=black, anchorcolor=red, urlcolor=blue}
\usepackage{longtable}
\usepackage{tabularx}
\usepackage{abstract}
\usepackage{appendix}
\usepackage{multicol}
\usepackage{bmpsize}
\usepackage[all]{hypcap}
\usepackage{titlesec}
\usepackage{indentfirst}
\usepackage{lipsum,titletoc}

%\setcounter{secnumdepth}{4}

%****************INIZIO GESTIONE SUBSECTION MULTIPLE
\makeatletter
\newcommand\level[1]{%
\ifcase#1\relax\expandafter\chapter\or
\expandafter\section\or
\expandafter\subsection\or
\expandafter\subsubsection\else
\def\next{\@level{#1}}\expandafter\next
\fi}
\newcommand{\@level}[1]{%
\@startsection{level#1}
{#1}
{\z@}%
{-3.25ex\@plus -1ex \@minus -.2ex}%
{1.5ex \@plus .2ex}%
{\normalfont\normalsize\bfseries}}

\newdimen\@leveldim
\newdimen\@dotsdim
{\normalfont\normalsize
\sbox\z@{0}\global\@leveldim=\wd\z@
\sbox\z@{.}\global\@dotsdim=\wd\z@
}

\newcounter{level4}[subsubsection]
\@namedef{thelevel4}{\thesubsubsection.\arabic{level4}}
\@namedef{level4mark}#1{}
\def\l@section{\@dottedtocline{1}{0pt}{\dimexpr\@leveldim*4+\@dotsdim*1+6pt\relax}}
\def\l@subsection{\@dottedtocline{2}{0pt}{\dimexpr\@leveldim*5+\@dotsdim*2+6pt\relax}}
\def\l@subsubsection{\@dottedtocline{3}{0pt}{\dimexpr\@leveldim*6+\@dotsdim*3+6pt\relax}}
\@namedef{l@level4}{\@dottedtocline{4}{0pt}{\dimexpr\@leveldim*7+\@dotsdim*4+6pt\relax}}

\count@=4
\def\@ncp#1{\number\numexpr\count@+#1\relax}
\loop\ifnum\count@<100
\begingroup\edef\x{\endgroup
\noexpand\newcounter{level\@ncp{1}}[level\number\count@]
\noexpand\@namedef{thelevel\@ncp{1}}{%
\noexpand\@nameuse{thelevel\@ncp{0}}.\noexpand\arabic{level\@ncp{1}}}
\noexpand\@namedef{level\@ncp{1}mark}####1{}%
\noexpand\@namedef{l@level\@ncp{1}}%
{\noexpand\@dottedtocline{\@ncp{1}}{0pt}{\the\dimexpr\@leveldim*\@ncp{5}+\@dotsdim*\@ncp{0}\relax}}}%
\x
\advance\count@\@ne
\repeat
\makeatother
\setcounter{secnumdepth}{100}
\setcounter{tocdepth}{100}
%****************FINE GESTIONE SUBSECTION MULTIPLE

%impostazioni relative alla visualizzazione delle section
%nell'indice
\titlecontents{section}
[0pt]%left indent
{\bfseries}
{\contentslabel{2.3em}}
{\hspace*{-2.3em}}
{\hfill\contentspage}
[]%separator


\oddsidemargin=.15in
\evensidemargin=.15in
\textwidth=6in
\topmargin=-.5in
\parindent=0in
\headheight=1in
\DeclareMathSizes{10}{10}{10}{10} %per piano qualifica
\pagestyle{fancy}
\lhead{
\bfseries {\Large \TipoDocumento}\\
\bfseries Versione: \Versione\\
}
\chead{}
\lhead{
\includegraphics[scale=0.455]{../Logo&Header/apertureHead.png}
}
\lfoot{\bfseries \TipoDocumento{} v\Versione}
\cfoot{}
\rfoot{\thepage\ of \mypageref{LastPage}}
\newcommand{\mypageref}[1]{
\hypersetup{linkcolor=black}\pageref{#1}\hypersetup{linkcolor=black}}
%\userpackage{lipsum}
\renewcommand{\footrulewidth}{0.4pt}

%definizioni comandi comuni utilizzati
\newcommand{\numref}[1]{\textsl{\nameref{#1} (\ref{#1})}}
\newcommand{\NomeGruppo}{Aperture Software}
\newcommand{\Progetto}{MaaP: MongoDB as an admin Platform}
\newcommand{\Prop}{CoffeeStrap}

%definizione tecnologie
\newcommand{\Node}{Node.js}
\newcommand{\NodeJS}{Node.js}
\newcommand{\Nodejs}{Node.js}

\newcommand{\mongodb}{MongoDB}

%tanti sub quanti ne vogliamo! :)
\newcommand{\subsubsubsection}{\level{4}}
\newcommand{\subsubsubsubsection}{\level{5}}
\newcommand{\subsubsubsubsubsection}{\level{6}}
\newcommand{\subsubsubsubsubsubsection}{\level{7}}
\newcommand{\subsubsubsubsubsubsubsection}{\level{8}}


%definizione comando per parola glossario
\newcommand{\gloss}[1]{\emph{#1}\ped{\emph{\tiny{G}}}}

\newcommand{\grassetto}{\textbf}

%per inserire immagini
\newcommand{\immagine}[2]{
\begin{center}
\begin{figure}[H]
\includegraphics[width=\textwidth]{{{#1}}}
\caption{#2}
\label{#1}
\end{figure}
\end{center}
}

\newcommand{\Glossario}{
Al fine di evitare ogni ambiguità nella comprensione del linguaggio utilizzato nel presente documento e, in generale, nella documentazione fornita dal gruppo \NomeGruppo{}, ogni termine tecnico, di difficile comprensione o di necessario approfondimento verrà inserito nel documento \emph{Glossario\_{}v\versioneGlossario{}.pdf}.\\
Saranno in esso definiti e descritti tutti i termini in corsivo e allo stesso tempo marcati da una lettera "G" maiuscola in pedice nella documentazione fornita.
}

\newcommand{\Prodotto}{
Lo scopo del prodotto è produrre un \gloss{framework} per generare interfacce web di amministrazione dei dati di business basato sullo stack \gloss{\Nodejs{}} e \gloss{\mongodb{}}.\\
L'obiettivo è quello di semplificare il lavoro allo \gloss{sviluppatore} che dovrà rispondere in modo rapido e standard alle richieste degli esperti di \gloss{business}.
}

%inizio pagina del documento
\begin{document}
\thispagestyle{empty}

\begin{center}\centerline{
%inserisco il logo grande della prima pagina
\includegraphics[scale=0.8]{../Logo&Header/logo.png}}

%metto il link dell'email sotto al logo
%{\href{mailto:ApertureSWE@gmail.com}{\color[rgb]{0.39,0.37,0.38}%ApertureSWE@gmail.com}}\\ [3pc]

\vspace{0.5in}

%titolo del progetto
{\Huge {\Progetto}}\\[.5pc]

\underline{\hspace{6in}}\\[8pc]

{\Huge {\TipoDocumento}}\\[1pc]
%{\emph{Versione \Versione}}\\
\end{center}

%\vspace{.05in}


%\vspace{.05in}

%informazioni documento
\begin{center}
%\section{Informazioni documento}
\begin{tabular}{r|l}
%\textbf{Nome} &\TipoDocumento \\
\textbf{Versione} & \Versione{} \\
\textbf{Data creazione} & \Data{} \\
\textbf{Data ultima modifica} & \DataUltimaModifica{} \\
\textbf{Stato del Documento} & Formale \\		%CAMBIARE QUI
\textbf{Uso del Documento} & Esterno \\			%CAMBIARE QUI
\textbf{Redazione} & Giacomo Pinato, Alessandro Benetti,\\			%CAMBIARE QUI
				   & Andrea Perin\\
\textbf{Verifica} & Michele Maso, Alberto Garbui\\				%ED ANCHE QUI!
\textbf{Approvazione} & Giacomo Pinato\\				%CAMBIARE QUI
\textbf{Distribuzione} & \parbox[t]{4cm}{Prof. Tullio Vardanega \\ Prof. Riccardo Cardin \\ \Prop{} }\\
\end{tabular}
\end{center}

\vspace{0.05in}

%inizio sommario del documento
\begin{abstract}
\begin{center}
Questo documento si propone di presentare la pianificazione del progetto MaaP.
\end{center}
\end{abstract}

%\vspace{.4in}

%seconda pagina, diario delle modifiche
\newpage
Diario delle modifiche
\begin{center}
\begin{longtable}{|c|c|c|p{0.5\linewidth}|}
\toprule
\textbf{Versione} & \textbf{Data} & \textbf{Autore} & \textbf{Modifiche effettuate}\\

%aggiungere qui una midrule per aggiungere una nuova riga alla tabella

\midrule
1.2.0 & 2013-12-18 & Giacomo Pinato (RE) & Approvazione documento\\
\midrule
1.1.1 & 2013-12-17 & Michele Maso (VE) & Verifica documento\\
\midrule
1.1.0 & 2013-12-16 & Alberto Garbui (VE) & Verifica documento\\
\midrule
1.0.9 & 2013-12-14 & Andrea Perin (RE) & Aggiunto consuntivo\\
\midrule
1.0.8 & 2013-12-06 & Alessandro Benetti (RE) & Aggiunto prospetto economico\\
\midrule
1.0.6 & 2013-12-05 & Alessandro Benetti (RE) & Aggiunto ciclo di vita\\
\midrule
1.0.5 & 2013-12-03 & Alessandro Benetti (RE) & Aggiunto analisi dei rischi\\
\midrule
1.0.4 & 2013-12-01 & Andrea Perin (RE) & Aggiunto gestione risorsa\\
\midrule
1.0.3 & 2013-12-01 & Giacomo Pinato (RE) & Aggiunto calendario scadenze\\
\midrule
1.0.2 & 2013-11-29 & Giacomo Pinato (RE) & Aggiunto organigramma\\
\midrule
1.0.1 & 2013-11-28 & Giacomo Pinato (RE) & Creazione documento\\

\bottomrule
\caption{Registro delle modifiche}
\label{tab:changelog}
\end{longtable}
\end{center}

%terza pagina Indice (viene aggiornato in automatico con due compilazioni)
\newpage
\tableofcontents

%pagine successive hanno la lista di tabelle e lista delle figure
%(vengono aggiornate in automatico)
\newpage
\listoftables
\listoffigures

%qui inizia la prima pagina ufficiale
\newpage
\section{Organigramma}%1.0
\label{1.0}
\subsection{Redazione}%1.1
\label{1.1}
\begin{center}
\begin{longtable}{|c|c|p{5cm}|}
\toprule
\textbf{Redattore} & \textbf{Data} & \textbf{Firma}\\
\midrule
Pinato Giacomo & 2013-11-28 & \\
\bottomrule
\end{longtable}
\end{center} 

\subsection{Approvazione}
\label{1.2}
\begin{center}
\begin{longtable}{|c|c|p{5cm}|}
\toprule
\textbf{Nome} & \textbf{Data} & \textbf{Firma}\\
\midrule
Pinato Giacomo & 2013-11-28 & \\
Vardanega Tullio &   & \\
\bottomrule
\end{longtable}
\end{center} 

\subsection{Accettazione componenti}
\label{1.3}
\begin{center}
\begin{longtable}{|c|c|p{5cm}|}
\toprule
\textbf{Nome} & \textbf{Data} & \textbf{Firma}\\
\midrule
Pinato Giacomo & 2013-11-26 & \\
Miotto Fabio & 2013-11-26 & \\
Maso Michele & 2013-11-26 & \\
Garbui Alberto & 2013-11-26 & \\
Sorgato Mattia & 2013-11-26 & \\
Perin Andrea & 2013-11-26 & \\
Benetti Alessandro & 2013-11-26 & \\
\bottomrule
\end{longtable}
\end{center} 

\subsection{Componenti}
\label{1.4}
\begin{center}
\begin{longtable}{|c|c|c|}
\toprule
\textbf{Nome} & \textbf{Matricola} & \textbf{Email}\\
\midrule
Pinato Giacomo & 1004030 & giacomo.pinato@gmail.com
\\
Miotto Fabio & 1003810 & fabietto.mi8@gmail.com\\
Maso Michele & 1004972 & maso.michele@gmail.com
\\
Garbui Alberto & 561226 & alberto.garbui@gmail.com\\
Sorgato Mattia & 1004404 & mattia.sorgato@gmail.com \\
Perin Andrea & 1037255 & a.xin90@gmail.com\\
Benetti Alessandro & 510890 & alexbenets@gmail.com\\
\bottomrule
\end{longtable}
\end{center} 

\subsection{Rotazione dei ruoli}
\label{1.5}
Durante lo sviluppo del progetto i ruoli che i membri del gruppo andranno a ricoprire saranno turnati a rotazione in modo che ogni membro abbia l'opportunità di adempire a tutti gli incarichi. Indispensabile è che non vi siano periodi in cui una stessa risorsa sia verificatrice di se stessa.

\newpage
\section{Ruoli di progetto}
\label{2.0}
Per la piena riuscita del progetto è indispensabile distinguere i vari ruoli che concorrono alla creazione del prodotto finale e le loro diverse responsabilità e competenze.
Ogni ruolo avrà una specifica area di competenza, degli specifici compiti e oneri e delle particolari autorizzazioni. Ogni componente dovrà limitarsi ai compiti ad esso assegnati e, nel caso qualcosa esuli dal suo campo di pertinenza, lo stesso dovrà rivolgersi al suo collega attualmente occupante il ruolo competente.
Tutti i membri, a rotazione, dovranno occupare come minimo una volta ciascuno dei ruoli descritti sottostante.

\subsection{Responsabile di Progetto}
\label{2.1}
Il Responsabile di Progetto incentra su di sé le responsabilità di scelta ed approvazione dei lavori. Ha inoltre il ruolo di rappresentare il gruppo nei contatti con l'esterno e durante la presentazione dei lavori.\\

Le sue competenze principali comprendono:
\begin{itemize}
\item Pianificazione, coordinamento e controllo delle attività;
\item Gestione e controllo delle risorse;
\item Approvazione delle analisi di gestione e rischio;
\item Approvazione dei documenti;
\item Comunicazioni con i committenti/proponenti.
\end{itemize}

Il responsabile ha il compito di assicurarsi che le attività di verifica vengano svolte sistematicamente seguendo le Norme di Progetto, che vengano rispettati i ruoli e le competenze assegnate nel Piano di Progetto e che non vi siano conflitti di interesse tra redattori e verificatori. Ha inoltre l'onere di gestire la creazione e l'assegnazione dei ticket di pianificazione e di assegnare ad un membro del gruppo il ruolo di responsabile di quest'ultimo, nel caso riguardi una sotto-attività.

\subsection{Amministratore}
\label{2.2}
L'Amministratore è il responsabile del controllo, dell'efficienza e dell'operatività dell'ambiente di lavoro e degli strumenti per la condivisone e la sincronizzazione.

\begin{itemize}
\item L'individuazione e la gestione di strumenti per automatizzare quanto più possibile processi o attività;
\item L'individuazione e la gestione di strumenti per il controllo dei processi e delle risorse;
\item L'individuazione e la gestione di strumenti e strategie per il controllo della qualità;
\item Gestione del versionamento.
\end{itemize}

\subsection{Analista}
\label{2.3}
L'analista è il responsabile dell'analisi dei requisiti di progetto. Dopo aver compreso pienamente la natura del problema e tutti i suoi domini, il suo ruolo è delineare vincoli e caratteristiche del prodotto finale, redigendo una specifica di progetto dettagliata, precisa e non ambigua, comprensibile sia dal proponente che dal progettista.

\subsection{Progettista}
\label{2.4}
Il progettista è colui che disegna una soluzione attuabile ed efficace che soddisfi i requisiti dettati dagli analisti. Il suo compito è progettare un'architettura che assicuri una facile manutenibilità del prodotto e una buona scomposizione in moduli indipendenti tra di loro.

\subsection{Verificatore}
\label{2.5}
Il Verificatore è responsabile delle attività di verifica. Ha il compito di assicurare che i  documenti e il codice rispettino gli standard qualitativi precedentemente definiti utilizzando gli strumenti e i metodi proposti dal Piano di Qualifica e attenendosi a quanto descritto nelle Norme di Progetto. 

\subsection{Programmatore}
\label{2.6}
Il Programmatore è responsabile delle attività di codifica e delle componenti di ausilio
necessarie per l'esecuzione delle prove di verifica e validazione.\\
Le responsabilità di  tale ruolo sono:
\begin{itemize}
\item Implementare rigorosamente le soluzioni descritte dal Progettista per la realizzazione del progetto;
\item Scrivere codice e la sua relativa documentazione che rispettino gli standard stabiliti per la loro scrittura;
\item Implementare i test sul codice scritto, necessari per prove di verifica e validazione.
\end{itemize}

\begin{center}
\begin{longtable}{|c|c|}
\toprule
\textbf{Ruolo} & \textbf{Costo in euro/ora}\\
\midrule
Responsabile & 30\\
Amministratore & 20\\
Analista & 25\\
Progettista & 22\\
Verificatore & 15\\
Programmatore & 15\\
\bottomrule
\caption{Tabella costi per ruolo}
\label{tab:costiruolo}
\end{longtable}
\end{center} 

\newpage
\section{Introduzione}
\label{3.0}
\subsection{Scopo del documento}
\label{3.1}
Il presente documento ha l'obiettivo di definire la pianificazione secondo la quale saranno scadenziati i lavori dal gruppo Aperture Software sul progetto \Progetto{}.\\

Gli scopi del presente documento sono:
\begin{itemize}
\item Presentare la pianificazione dei tempi e delle attività;
\item Preventivare l'utilizzo delle risorse;
\item Consuntivare l'utilizzo delle risorse durante l'evolversi dei lavori;
\item Analizzare i possibili fattori di rischio.
\end{itemize}

\subsection{Scopo del prodotto}
\label{3.2}
\Prodotto{}

\subsection{Glossario}
\label{3.3}
\Glossario{}

\subsection{Riferimenti}
\label{3.4}
\subsubsection{Normativi}
\label{3.4.1}
\begin{itemize}
\item \grassetto{Capitolato d'appalto C1}:\\ \url{http://www.math.unipd.it/~tullio/IS-1/2013/Progetto/C1.pdf}
\item Norme di Progetto: Norme\_{}di\_{}Progetto\_{}v\versioneNormeDiProgetto{}.pdf  (allegato alla presente documentazione)\\
\end{itemize}

\subsubsection{Informativi}
\label{3.4.2}
\begin{itemize}
\item \grassetto{Software Engineering} - Ian Sommerville - 9th Edition (2010): 
– Part 4: Software Management.
\item \grassetto{Slide dell'insegnamento Ingegneria del Software modulo A}:\\
– Il ciclo di vita del software;
– Gestione di progetto.\\
\url{http://www.math.unipd.it/~tullio/IS-1/2013/};
\item \grassetto{Metriche di progetto}:\\
\url{http://it.wikipedia.org/wiki/Metriche_di_progetto}.
\end{itemize}

\subsection{Ciclo di vita}
\label{3.5}
Un processo software è costituito da un insieme di attività che conducono alla realizzazione di un prodotto software. Abbiamo deciso di utilizzare il modello incrementale perché grazie alle sue caratteristiche lo riteniamo più idoneo e flessibile per fare un lavoro di gruppo:
\begin{itemize}
\item Il problema viene scomposto in più sotto-problemi in modo da permettere una suddivisione più semplice delle risorse e dei tempi in quanto, tutte le risorse vengono utilizzare per un numero limitato di attività in un breve lasso di tempo. Questo rende più facile l'esecuzione dei test perché essi saranno più dettagliati e quindi più esaustivi;
\item I requisiti del progetto vengono trattati in base ad un ordine di priorità ottenuto ordinando i requisiti per la loro importanza strategica e quindi vengono svolti per primi quelli di maggiore criticità;
\item Il rischio di fallimento viene ridotto in quanto ogni incremento consolida soltanto la sezione coinvolta;
\item I rilasci del software sono multipli e successivi in quanto nei primi rilasci saranno relativi ai requisiti di primaria importanza, in modo che tali requisiti subiscano il maggior numero attività di verifica e risultino quindi più raffinati e migliorati, mentre negli ultimi rilasci si andranno ad aumentare il numero di funzionalità del prodotto e un miglioramento delle funzionalità già in essere.
\end{itemize}

Il modello quindi permette che il committente in corso d'opera possa già vedere dei prototipi con le funzionalità di primaria importanza già implementate e possa quindi valutare il lavoro in corso d'opera e suggerire eventuali correzioni e miglioramenti.


\subsection{Scadenze}
\label{3.6}
\NomeGruppo{} ha deciso di rispettare le seguenti scadenze per la pianificazione del progetto:
\begin{itemize}
\item Revisione dei Requisiti: 2013-12-20;
\item Revisione di Progetto:  2014-03-18;
\item Revisione di Qualifica: 2014-07-07;
\item Revisione di Accettazione: 2014-07-21.
\end{itemize}

\newpage
\section{Pianificazione}
\label{4.0}
Avendo scelto di operare con il ciclo incrementale e per rispettare le scadenze come esposte nel punto precedente, abbiamo deciso di dividere il progetto in cinque fasi incrementali distinte.

\begin{enumerate}
\item Analisi (AR);
\item Analisi di Dettaglio (AD);
\item Progettazione Architetturale (PA);
\item Progettazione di Dettaglio e Codifica (PC);
\item Verifica e Validazione (VV).
\end{enumerate}

Ognuna delle fasi del progetto è stata poi divisa in attività in modo da poter associare ad ogni attività una o più risorse. Per avere una miglior organizzazione le attività sono state poi spezzante in sotto-attività. 

\subsection{Analisi}


\subsubsection{Diagramma di Gantt delle attività}
\immagine{analisidettaglio}{Diagramma di Gantt, fase di Analisi}

\subsubsection{Diagramma WBS delle attività}
\immagine{analisidettaglioWBS}{Work Breakdown Structure, fase di Analisi}

\newpage
\subsection{Analisi Dettaglio}
Questa attività ha inizio in data 2013-12-23 e avrà fine e avrà fine in data 2014-01-04.\\
Questa attività inizia dopo la \grassetto{Revisione dei Requisiti} e termina con l'inizio della attività di \grassetto{Progettazione Architetturale}.

\subsubsection{Diagramma di Gantt delle attività}
\immagine{./gantt/analisidettaglio}{Diagramma di Gantt, attività di Analisi Dettaglio}

\newpage
\subsubsection{Diagramma WBS delle attività}
\immagine{./gantt/analisidettaglioWBS}{Work Breakdown Structure, attività di Analisi Dettaglio}

\newpage
\subsubsection{Ripartizione ore}
\immagine{./tabelle/allocazione_dettaglio}{Allocazione risorse, attività di Analisi Dettaglio}

\newpage
Nell'attività di \grassetto{Analisi Dettaglio}, ciascun componente dovrà rivestire i seguenti ruoli:
\immagine{./tabelle/ore_dettaglio}{Ore per componente, attività di Analisi Dettaglio}

Il seguente grafico mostra la suddivisione dei ruoli e il monte ore ricoperto da ciascun componente del gruppo:
\immagine{./grafici/orecomponente_dettaglio}{Ore per componente, attività di Analisi Dettaglio}


\newpage
\subsection{Progettazione Architetturale}
Questa attività ha inizio in data 2014-01-13 e avrà fine in data 2014-03-18 per un totale di 65 giorni. In questa attività sono attivi i ruoli di Amministratore, Analista, Progettista, Responsabile e Verificatore.

\subsubsection{Diagramma di Gantt delle attività}
\immagine{./gantt/progett}{Diagramma di Gantt, attività di Progettazione Architetturale}

\newpage
\subsubsection{Diagramma WBS delle attività}
\immagine{./gantt/progettWBS}{Work Breakdown Structure, attività di Progettazione Architetturale}

\newpage
\subsubsection{Ripartizione ore}
\immagine{./tabelle/allocazione_proge}{Allocazione risorse, attività di Progettazione Architetturale}

\newpage
Nell'attività di \grassetto{Progettazione Architetturale}, ciascun componente dovrà rivestire i seguenti ruoli:
\immagine{./tabelle/ore_proge}{Ore per componente, attività di Progettazione Architetturale}

Il seguente grafico mostra la suddivisione dei ruoli e il monte ore ricoperto da ciascun componente del gruppo:
\immagine{./grafici/orecomponente_proge}{Ore per componente, attività di Progettazione Architetturale}


\newpage
\subsection{Progettazione di Dettaglio e Codifica}
Questa fase ha inizio in data 2014-03-24 e avrà fine in data 2014-07-07 per un totale di 105 giorni. In questa fase sono attivi i ruoli di Amministratore, Analista (per un ultimissimo ritocco all'Analisi dei requisiti, se necessario), Progettista, Responsabile, Programmatore e Verificatore.

\subsubsection{Diagramma di Gantt delle attività}
\immagine{./gantt/progcodifica}{Diagramma di Gantt, fase di Progettazione di Dettaglio e Codifica}

\newpage
\subsubsection{Diagramma WBS delle attività}
\immagine{./gantt/progcodificaWBS}{Work Breakdown Structure, fase di Progettazione di Dettaglio e Codifica}

\newpage
\subsubsection{Ripartizione ore}
\immagine{./tabelle/allocazione_codifica}{Allocazione risorse, fase di Progettazione di Dettaglio e Codifica}

\newpage
Nella fase di \grassetto{Progettazione di Dettaglio e Codifica}, ciascun componente dovrà rivestire i seguenti ruoli:
\immagine{./tabelle/ore_codifica}{Ore per componente, fase di Progettazione di Dettaglio e Codifica}

Il seguente grafico visualizza la suddivisione dei ruoli e per quanto essi sono stati ricoperti da ciascun membro del gruppo:
\immagine{./grafici/orecomponente_codifica}{Ore per componente, fase di Progettazione di Dettaglio e Codifica}

\newpage
\subsection{Verifica e Validazione}


\subsubsection{Diagramma di Gantt delle attività}
\immagine{./gantt/verifica}{Diagramma di Gantt, fase di Verifica e Validazione}

\newpage
\subsubsection{Diagramma WBS delle attività}
\immagine{./gantt/verificaWBS}{Work Breakdown Structure, fase di Verifica e Validazione}

\newpage
\subsubsection{Ripartizione ore}
\immagine{./tabelle/allocazione_verifica}{Allocazione risorse, fase di Verifica e Validazione}

\newpage
Nella fase di \grassetto{Verifica e Validazione}, ciascun componente dovrà rivestire i seguenti ruoli:
\immagine{./tabelle/ore_verifica}{Ore per componente, fase di Verifica e Validazione}

Il seguente grafico visualizza la suddivisione dei ruoli e per quanto essi sono stati ricoperti da ciascun membro del gruppo:
\immagine{./grafici/orecomponente_verifica}{Ore per componente, fase di Verifica e Validazione}

\newpage
\subsection{Prospetto Economico}

\subsubsection{Analisi}

\immagine{./tabelle/costo_analisi}{Costo per ruolo, fase di Analisi}

\immagine{./grafici/oreruolo_analisi}{Ore per ruoli, fase di Analisi}

\immagine{./grafici/costoruolo_analisi}{Costo per ruoli, fase di Analisi}


\subsubsection{Analisi Dettaglio}

\immagine{./tabelle/costo_dettaglio}{Costo per ruolo, fase di Analisi Dettaglio}

\immagine{./grafici/oreruolo_dettaglio}{Ore per ruoli, fase di Analisi Dettaglio}

\immagine{./grafici/costoruolo_dettaglio}{Costo per ruoli, fase di Analisi Dettaglio}

\subsubsection{Progettazione Architetturale}

\immagine{./tabelle/costo_proge}{Costo per ruolo, fase di Progettazione Architetturale}

\immagine{./grafici/oreruolo_proge}{Ore per ruoli, fase di Progettazione Architetturale}

\immagine{./grafici/costoruolo_proge}{Costo per ruoli, fase di Progettazione Architetturale}

\subsubsection{Progettazione di Dettaglio e Codifica}

\immagine{./tabelle/costo_codifica}{Costo per ruolo, fase di Progettazione di Dettaglio e Codifica}

\immagine{./grafici/oreruolo_codifica}{Ore per ruoli, fase di Progettazione di Dettaglio e Codifica}

\immagine{./grafici/costoruolo_codifica}{Costo per ruoli, fase di Progettazione di Dettaglio e Codifica}

\subsubsection{Verifica e Validazione}

\immagine{./tabelle/costo_verifica}{Costo per ruolo, fase di Verifica e Validazione}

\immagine{./grafici/oreruolo_verifica}{Ore per ruoli, fase di Verifica e Validazione}

\immagine{./grafici/costoruolo_verifica}{Costo per ruoli, fase di Verifica e Validazione}


\subsubsection{Totale}
Nella tabella seguente sono riportate le ore totali di investimento previste per la realizzazione dell'intero progetto, assieme alle ore rendicontate:
\immagine{./tabelle/ore_totali}{Costi totali per ruolo}

I seguenti grafici descrivono quanto ogni ruolo abbia partecipato rispettivamente in ore complessive ed ore retribuite alla realizzazione del progetto: 
\immagine{./grafici/oreruolo_totali}{Ore totali per ruoli}
\immagine{./grafici/oreruolo_retribuite}{Ore totali retribuite per ruoli}

Infine i seguenti diagrammi descrivono il carico di lavoro assegnato ad ogni componente del gruppo per ogni ruolo rispettivamente per le ore totali del progetto e per le ore retribuite:
\immagine{./grafici/orecomponente_totali}{Ore totali per componente}
\immagine{./grafici/orecomponente_retribuite}{Ore retribuite per componente}

\subsubsection{Conclusioni}
In conclusione le ore complessive delle fasi di Progettazione Architetturale, Progettazione di Dettaglio e Codifica, Verifica e Validazione ammontano a 733 ore con un costo preventivato di euro \grassetto{13801} che sarà interamente a carico del Proponente. 

%parte finale, analisi dei rischi
\newpage
\section{Analisi dei rischi}
Nello sviluppo del progetto si è dedicato parte del tempo per individuare rischi che potrebbero compromettere la realizzazione del progetto.\\
Per prevedere e controllare eventuali rischi si è seguita questa sequenza di passi:
\begin{itemize}
\item \grassetto{Identificazione e analisi:} si cerca di individuare i rischi che possono comparire durante lo sviluppo del progetto e nella realizzazione del prodotto; successivamente con l'analisi si cerca di comprendere eventuali criticità e conseguenze che questi rischi possono portare;
\item \grassetto{Controllo:} si valutano le strategie da attuare  per prevenire i rischi;
\item \grassetto{Soluzione:} azioni da intraprendere, in caso di rischio già avvenuto, per vanificare l'effetto dello stesso.
\item \grassetto{Riscontro:} vengono riportati i riscontri effettivi verificati nel corso del progetto. In questo punto verranno anche riassunte le conclusioni che questa verifica ha portato e le attuazioni sul rischio stesso. 
\end{itemize}
Di seguito verranno elencati i rischi emersi e saranno accompagnati da una breve descrizione, un livello di rischio, le contromisure da adottare e l'impatto che essi avranno. I rischi sono stati ordinati per probabilità di occorrenza e impatto dal più alto al più basso.\\
I livelli di rischio e gli impatti di ogni singolo rischio sono posizionati su una scala di 5 valori:
\begin{enumerate}
\item \grassetto{alto;}
\item \grassetto{medio-alto;}
\item \grassetto{medio;}
\item \grassetto{medio-basso;}
\item \grassetto{basso.}
\end{enumerate}

\subsection{Problemi componenti del gruppo}
\begin{itemize}
\item \grassetto{Descrizione}: ogni componente del gruppo ha impegni personali e svolge attività extra accademiche ed alcuni di essi svolgono un vero e proprio lavoro. Di conseguenza non possono essere spesso disponibili;
\item \grassetto{Livello di rischio}: alto;
\item \grassetto{Contromisure}: è compito del Responsabile di Progetto stilare appositi calendari di gruppo per organizzare le giornate lavorative per il progetto, in modo tale da non contrastare gli impegni personali di ciascun componente. Inoltre il carico di lavoro che un componente non svolgerà per impegni personali, dovrà essere ripartito tra gli altri componenti del gruppo;
\item \grassetto{Impatto}: alto.
\item \grassetto{Riscontro}: Sebbene la maggior parte del gruppo si ritrovi sempre negli orari stabiliti, alcuni membri, dati li impegni personali, non sono sempre operativi. Questo rischio è stato riscontrato con alta occorrenza, e ha portato ad un rallentamento nell'avanzamento del progetto. Quindi i livelli di rischio e impatto rimangono inalterati. Sarà compito del Responsabile di Progetto stilare dei piani più stringenti o prendere provvedimenti individuali.
\end{itemize}

\subsection{Conoscenze tecnologiche}
\begin{itemize}
\item \grassetto{Descrizione}: i componenti del gruppo non hanno una buona conoscenza delle tecnologie utilizzate, in quanto queste sono delle novità assolute per alcuni di loro, in particolare:
\begin{itemize}
\item \gloss{MongoDB};
\item \gloss{Node.js};
\item Mongoose;
\item \gloss{Express}.
\end{itemize}
La scarsa conoscenza delle precedenti tecnologie saranno un ostacolo che i componenti del gruppo troveranno durante lo sviluppo del progetto.
\item \grassetto{Livello di rischio}: alto;
\item \grassetto{Contromisure}: i componenti del gruppo si impegnano per istruirsi sull'utilizzo delle tecnologie richieste, partecipando a seminari specifici e utilizzando materiale e documentazione reperibile sul \gloss{Web} o che fornisce l'Amministratore. Questi apprendimenti saranno fatti in tempi brevi, per non compromettere e ritardare il proseguimento del progetto; i componenti del gruppo hanno dedicato delle ore pianificate il fine settimana per studiare queste tecnologie necessarie;
\item \grassetto{Impatto}: medio. I componenti del gruppo incontreranno difficoltà nell'utilizzo delle tecnologie apprese.
\item \grassetto{Riscontro}: si è riscontrato che la quantità di materiale informativo/educativo riguardante le tecnologie necessarie è davvero molto alta. La quantità e la qualità di tale materiale ha portato ad una facile educazione dei membri del team, con basso dispendio di tempo ed energie rispetto a quanto previsto. Da queste premesse, è stato deciso di abbassare il livello di impatto da "alto" a "medio". 
\end{itemize}

\subsection{Cambi nelle tecnologie}
\begin{itemize}
\item \grassetto{Descrizione}: data la recentezza delle tecnologie adottate nello sviluppo di questo progetto, è possibile che le specifiche dello stack tecnologico utilizzato varino nel tempo. Questo probabilmente costringerebbe il gruppo a rivedere parte della codifica;
\item \grassetto{Livello di rischio}: medio;
\item \grassetto{Contromisure}: il gruppo dovrà tenersi costantemente informato sugli aggiornamenti delle versioni dei software. In caso nuove release, verranno effettuati dei test per garantire la compatibilità del codice prodotto;
\item \grassetto{Impatto}: medio-alto;
\item \grassetto{Riscontro}: da verificare.
\end{itemize}

\subsection{Comprensione requisiti}
\begin{itemize}
\item \grassetto{Descrizione}: è possibile che i componenti del gruppo non comprendano in pieno i requisiti e che alcuni aspetti vengano trattati in maniera errata o incompleta;
\item \grassetto{Livello di rischio}: medio;
\item \grassetto{Contromisure}: per ridurre al minimo gli effetti che il rischio comporta, ci saranno, durante il periodo di tempo dedicato all'Analisi dei Requisiti, degli incontri con il Proponente per delle delucidazioni in merito ai requisiti richiesti dal prodotto;
\item \grassetto{Impatto}: medio-alto. Se si verifica tale rischio è necessario aggiornare con piccole modifiche il documento di Analisi dei Requisiti, o nel caso peggiore sarà necessario stilare una nuova \gloss{versione} del medesimo documento.
\item \grassetto{Riscontro}: durante il progetto si sono verificati casi di occorrenza di questo rischio. In particolare, durante le attività di progettazione, la comprensione dei requisiti è stata a tratti ambigua e ciò ha portato a diverse consultazioni con gli Analisti e a revisioni della progettazione, anche se non particolarmente gravi. L'impatto sulla progettazione non è stato grave come previsto, ma comunque abbastanza alto. Si è deciso quindi di abbassare di mezzo grado l'impatto, da "alto" a "medio-alto".
\end{itemize}

\subsection{Preventivi di costi errati}
\begin{itemize}
\item \grassetto{Descrizione}: i componenti del gruppo, non avendo esperienza in ambito manageriale e finanziario, possono imbattersi in errori di pianificazione dei tempi e conseguente aumento dei costi;
\item \grassetto{Livello di rischio}: medio-basso;
\item \grassetto{Contromisure}: si cercherà di dedicare più tempo alla pianificazione dei tempi e costi, affrontando i vari passi che ne derivano con maggior calma e attenzione;
\item \grassetto{Impatto}: medio-alto. Se il rischio si verifica, bisognerà avvisare il Proponente e cercare di ridurre i costi per avvicinarsi il più possibile alla soglia prestabilita.
\item \grassetto{Riscontro}: da quanto rilevato dai consuntivi stilati, i preventivi non si sono discostati dai consuntivi in maniera rilevante. Si è deciso quindi di abbassare il livello di rischio da "medio" a "medio-basso".
\end{itemize}

\subsection{Variabilità requisiti}
\begin{itemize}
\item \grassetto{Descrizione}: il gruppo non deve sottovalutare la possibilità che il Proponente possa cambiare i requisiti in corso d'opera;
\item \grassetto{Livello di rischio}: basso;
\item \grassetto{Contromisure}: se il rischio si verifica, sarà compito del gruppo adeguarsi ai nuovi requisiti imposti dal Proponente;
\item \grassetto{Impatto}: medio-alto. Un cambiamento sostanziale dei requisiti porterà ad un impatto alto di tale rischio e al conseguente ritardo nella consegna del prodotto finale.
\item \grassetto{Riscontro}: il Proponente è stato sempre molto disponibile e flessibile negli incontri fissati, e in nessun caso ha mai avanzato richieste che esulassero da quelle descritte nel capitolato. Le uniche modifiche richieste sono state molto specifiche e comunque in linea con i nostri requisiti. Si è quindi deciso di abbassare il livello di rischio da "medio" a "basso".
\end{itemize}

\subsection{Problemi hardware}
\begin{itemize}
\item \grassetto{Descrizione}: gran parte del lavoro è basato sull'utilizzo di personal computer e di un \gloss{server} per il \gloss{repository} e per il sistema di gestione dei \gloss{ticket}, descritto nella sezione apposita delle Norme di Progetto. In seguito ad una loro eventuale rottura si potrebbe perdere parte del lavoro svolto e le conseguenti attività di ripristino porterebbero ad un ritardo nello svolgimento delle attività ripartite per il prosieguo del progetto;
\item \grassetto{Livello di rischio}: basso;
\item \grassetto{Contromisure}: sulla componente server verranno effettuati periodici backup per non perdere il lavoro svolto, ed anche i singoli componenti del gruppo, al termine della giornata lavorativa, effettueranno un backup su sistemi di \gloss{Cloud} o dispositivi \gloss{hardware} esterni;
\item \grassetto{Impatto}: basso. In caso di rischio avvenuto, la presenza dei numerosi backup faciliterà il compito dell' Amministratore, ovvero il ripristino del lavoro svolto.
\item \grassetto{Riscontro}: in base a quanto premesso, si è notato in effetti che i problemi hardware in genere si sono presentati con scarsa probabilità e non hanno influito sulla condotta del progetto. In questi casi hanno aiutato le strutture messe a disposizione dalla facoltà, le quali ci hanno permesso di lavorare con pari efficacia. Ci è parso logico mantenere i livelli di rischio e impatto inalterati.
\end{itemize}

\newpage
\section{Consuntivo e Preventivo a finire}
Questa sezione contiene il prospetto economico che riporta le spese effettivamente sostenute. Vengono riportate le ore impiegate per svolgere i compiti preventivati. In base alla differenza di ore tra il preventivo ed il consuntivo, detta conguaglio, avremmo un bilancio:
\begin{itemize}
\item \grassetto{Positivo}: se il preventivo ha superato il consuntivo;
\item \grassetto{Negativo}: se il consuntivo ha superato il preventivo;
\item \grassetto{In pari}: se preventivo e consuntivo coincidono.
\end{itemize}
\subsection{Analisi}
Di seguito è riportato il consuntivo del periodo di tempo dedicato all'Analisi dei Requisiti.\\
La tabella sottostante riporta la differenza delle ore tra preventivo e consuntivo, divise per ruolo.

\immagine{./tabelle/differenza_ruoli}{Costo per ruolo, fase di Pre RR}

Nella tabella seguente sono riportate le differenze tra le ore di lavoro previste per ogni componente con quelle realmente impiegate.

\immagine{./tabelle/differenza_preventivo_consuntivo}{Differenza preventivo consuntivo per componente, fase di Pre RR}

\subsubsection{Conclusioni}
In conclusione, come si può notare dai valori presenti nelle precedenti tabelle, è stata impiegata un'ora in più per svolgere le attività programmate con un bilancio in negativo di euro 35. Questa somma non andrà ad intaccare il costo totale del progetto, in quanto le ore spese in questo periodo non sono rendicontate e quindi non sono a carico del Proponente.

\subsection{Analisi in Dettaglio}
Di seguito è riportato il consuntivo del periodo di tempo dedicato all' Analisi in Dettaglio.\\
La tabella sottostante riporta la differenza delle ore tra preventivo e consuntivo, divise per ruolo.

\immagine{./tabelle/differenza_ruoli_analisidettaglio}{Costo per ruolo, fase di Post RR}

Nella tabella seguente sono riportate le differenze tra le ore di lavoro previste per ogni componente con quelle realmente impiegate.

\immagine{./tabelle/differenza_preventivo_consuntivo_analisidettaglio}{Differenza preventivo consuntivo per componente, fase di Post RR}

\subsubsection{Conclusioni}
L'Analisi dei Requisiti è stata corretta secondo le osservazioni fatte dal Committente, inoltre è stata incrementata dettagliandola maggiormente.\\
In conclusione, come si può notare dai valori presenti nelle precedenti tabelle, sono state impiegate due ore in più per svolgere le attività di Analisi programmate, con un bilancio in negativo di euro 50. Questa somma non andrà ad intaccare il costo totale del progetto, in quanto le ore spese in questo periodo non sono rendicontate e quindi non sono a carico del Proponente.

\subsection{Progettazione Architetturale}
Di seguito è riportato il consuntivo del periodo di tempo dedicato alla Progettazione Architetturale.\\
La tabella sottostante riporta la differenza delle ore tra preventivo e consuntivo, divise per ruolo.

\immagine{./tabelle/differenza_ruoli_progettaz}{Costo per ruolo, fase di Pre RP}

Nella tabella seguente sono riportate le differenze tra le ore di lavoro previste per ogni componente con quelle realmente impiegate.

\immagine{./tabelle/differenza_preventivo_consuntivo_progettaz}{Differenza preventivo consuntivo per componente, fase di Pre RP}

\subsubsection{Conclusioni}
Si è deciso di spostare la data per la revisione di progettazione dal 2014-03-18 al 2014-03-29, perchè, il tempo preventivato dedicato allo studio personale di tecnologie e di conoscenze specifiche per la progettazione, è stato molto sottostimato rispetto al tempo realmente occorso. Il tutto ha portato allo slittamento della consegna, ma il costo a budget speso per il tempo dedicato allo studio personale non è a carico del Proponente.
In conclusione, come si può notare dai valori presenti nelle precedenti tabelle, sono state tolte delle ore di Amministratore e di Analista, mentre sono state richieste delle ore di Verificatore in più rispetto a quanto pianificato; questa modifica ha portato ad un bilancio in positivo di 63 euro.
\subsection{Preventivo a finire}
Come si può notare da quanto scritto in precedenza, c'è stato un attivo pari a 63 euro; questa somma verrà sicuramente utilizzata per aumentare le ore di Verifica nel periodo di tempo che porta alla successiva revisione.\\
Si è deciso di investire questa somma nell'attività di Verifica nella fase di Pre RQ, così da migliorare la qualità dei prodotti in uscita. 


%FINE DOCUMENTO NON CANCELLARE
\end{document}